\section{Discussion}

For the most part, the general shape of the plots matched that of the simulations, with the largest exceptions being state one for the disturbance rejection PID and state one for reference tracking PID.
It is evident for state one for the disturbance rejection PID that the system was not cooled down to room temperature before the controlling commenced.
For state one for the reference tracking PID, various results were found, which are believed to be due to the already weak disturbance rejection causing a significant effect due to external disturbances.

The general PID controller in the Laplace domain is given by:

\begin{equation}
    C(s) = K_p + \frac{K_i}{s} + K_d s
    \label{eq:pid}
\end{equation}

Each term contributes differently to the closed-loop response, and the simulation results illustrate these effects clearly.

\subsection{Proportional Term}

The proportional term $u_P(t) = K_p\,e(t)$ produces a control action proportional to the current error.
It increases response speed but cannot eliminate steady-state error for a step input.
The steady-state output of a P-only closed-loop system is:

\begin{equation}
    y_{ss} = \frac{K_p\,G_{11}(0)}{1 + K_p\,G_{11}(0)} = \frac{0.8475 \times 15.59}{1 + 0.8475 \times 15.59} = 0.9296
    \label{eq:p_ss}
\end{equation}

This matches the simulated result in \cref{tab:p_metrics}, where the closed-loop output settles at $92.96\%$ of the \SI{50}{\degreeCelsius} reference, leaving a permanent error of approximately \SI{3.5}{\degreeCelsius} as seen in \cref{fig:p_nodist}.
In the disturbance case (\cref{fig:p_dist}), the P controller partially compensates by reducing $u(t)$, but a permanent offset remains since there is no mechanism to accumulate and correct the residual error.

\subsection{Integral Term}

The integral term $u_I(t) = K_i\int_0^t e(\tau)\,d\tau$ accumulates error over time, guaranteeing zero steady-state error for a constant reference by the Internal Model Principle.
This is confirmed in \cref{tab:pi_metrics}, where the PI controller achieves a steady-state gain of exactly $1.0000$, and \cref{fig:pi_nodist} shows $y(t)$ converging to exactly \SI{50}{\degreeCelsius}.

The cost is a slower, more oscillatory response: rise time increases from \SI{7.9}{\second} (P) to \SI{103.1}{\second} (PI), and overshoot increases from $7.42\%$ to $12.95\%$.
The integral term also fully rejects the constant disturbance, as seen in \cref{fig:pi_dist} where $y(t)$ returns to \SI{50}{\degreeCelsius} after the PWM2 disturbance at $t = \SI{350}{\second}$, unlike the P controller in \cref{fig:p_dist}.

\subsection{Derivative Term}

The derivative term $u_D(t) = K_d\,\frac{de}{dt}$ acts on the rate of change of error, providing a damping action that reduces overshoot.
The PD controller results in \cref{tab:pd_metrics} show an overshoot of only $8.8 \times 10^{-5}\%$, compared to $7.42\%$ for P alone, at the cost of a slower rise time (\SI{77.4}{\second}).
Like the P controller, PD retains a steady-state error since it contains no integrator.

\subsection{Controller Comparison}

Combining all three terms allows the PID controller to benefit from the speed of P, the accuracy of I, and the damping of D simultaneously.
The reference-tracking and disturbance-rejection PID controllers (\cref{subsec:pid_ref,subsec:pid_dist}) represent different tuning priorities within this structure: the reference-tracking design favours a well-damped step response, while the disturbance-rejection design increases $K_i$ and $K_d$ to suppress disturbances more aggressively.
